\expandafter\gdef\csname SNside1\endcsname{\node[anchor=north west,inner sep=0pt,text width=48mm,font=\SideFont,align=left] (SNsidebox) at ([xshift=138mm,yshift=-343.27pt]current page.north west) {\textcolor{SNink}{\textbf{课程连接 · 逻辑与物理}}\par\vspace{1.5mm}同一逻辑关系可以用不同存储结构实现。顺序或链式描述物理安排，线性或非线性描述元素之间的关系。};\pgfextracty{\SNsideheight}{\pgfpointdiff{\pgfpointanchor{SNsidebox}{south}}{\pgfpointanchor{SNsidebox}{north}}}\typeout{SNHEIGHT:1:14:\the\SNsideheight}}
\expandafter\gdef\csname SNside2\endcsname{\node[anchor=north west,inner sep=0pt,text width=48mm,font=\SideFont,align=left] (SNsidebox) at ([xshift=138mm,yshift=-95.00pt]current page.north west) {\textcolor{SNink}{\textbf{课程连接 · 逻辑与物理}}\par\vspace{1.5mm}同一逻辑关系可以用不同存储结构实现。顺序或链式描述物理安排，线性或非线性描述元素之间的关系。};\pgfextracty{\SNsideheight}{\pgfpointdiff{\pgfpointanchor{SNsidebox}{south}}{\pgfpointanchor{SNsidebox}{north}}}\typeout{SNHEIGHT:2:14:\the\SNsideheight}}
\expandafter\gdef\csname SNside3\endcsname{\node[anchor=north west,inner sep=0pt,text width=48mm,font=\SideFont,align=left] (SNsidebox) at ([xshift=138mm,yshift=-95.00pt]current page.north west) {\textcolor{SNink}{\textbf{错项排除 · 有穷不等于很快}}\par\vspace{1.5mm}执行2的n次方步，只要对每个合法有限n都能结束，仍可满足有穷性。可行性也不是``考试规模下跑得快''；先判基本特性，再评价实际效率，不要混用两个层次。};\pgfextracty{\SNsideheight}{\pgfpointdiff{\pgfpointanchor{SNsidebox}{south}}{\pgfpointanchor{SNsidebox}{north}}}\typeout{SNHEIGHT:3:2:\the\SNsideheight}}
\expandafter\gdef\csname SNside4\endcsname{\node[anchor=north west,inner sep=0pt,text width=48mm,font=\SideFont,align=left] (SNsidebox) at ([xshift=138mm,yshift=-95.00pt]current page.north west) {\textcolor{SNink}{\textbf{不变量 · 累计值与轮数不同}}\par\vspace{1.5mm}第k轮结束时sum=k(k+1)/2，循环次数是满足停止不等式的最小k。sum约等于n不代表循环n次；每轮增量逐步变大，正是平方根阶出现的原因。};\pgfextracty{\SNsideheight}{\pgfpointdiff{\pgfpointanchor{SNsidebox}{south}}{\pgfpointanchor{SNsidebox}{north}}}\typeout{SNHEIGHT:4:6:\the\SNsideheight}}
\expandafter\gdef\csname SNside5\endcsname{\node[anchor=north west,inner sep=0pt,text width=48mm,font=\SideFont,align=left] (SNsidebox) at ([xshift=138mm,yshift=-95.00pt]current page.north west) {\textcolor{SNink}{\textbf{不变量 · 先列每轮工作量}}\par\vspace{1.5mm}外层i取1、2、4等值，内层工作量恰与i相同。用等比和可知所有小轮之和仍受末轮的常数倍控制；``两层循环''没有提供平方阶的充分条件。};\pgfextracty{\SNsideheight}{\pgfpointdiff{\pgfpointanchor{SNsidebox}{south}}{\pgfpointanchor{SNsidebox}{north}}}\typeout{SNHEIGHT:5:7:\the\SNsideheight}}
\expandafter\gdef\csname SNside6\endcsname{\node[anchor=north west,inner sep=0pt,text width=48mm,font=\SideFont,align=left] (SNsidebox) at ([xshift=138mm,yshift=-95.00pt]current page.north west) {\textcolor{SNink}{\textbf{最小反例 · 链式不等于树}}\par\vspace{1.5mm}三个结点a\ensuremath{\to}b\ensuremath{\to}c\ensuremath{\to}a可以构成循环单链表的物理链接，链接图含环，不能据此认定为树。采用指针联系只是存储表示；逻辑对象究竟是线性表、树还是图，还需查看关系与操作约定。\par\smallskip\begin{tikzpicture}[x=1mm,y=1mm,>=Stealth,node distance=5mm]
\node[draw=SNink!50,fill=SNblue,rounded corners=1mm,text width=34mm,align=center,inner sep=1.2mm,font=\fontsize{7.5}{9}\selectfont] (m0) at (0,0) {a};
\node[draw=SNink!50,fill=SNblue,rounded corners=1mm,text width=34mm,align=center,inner sep=1.2mm,font=\fontsize{7.5}{9}\selectfont] (m1) at (0,-11) {b};
\node[draw=SNink!50,fill=SNblue,rounded corners=1mm,text width=34mm,align=center,inner sep=1.2mm,font=\fontsize{7.5}{9}\selectfont] (m2) at (0,-22) {c};
\draw[->,SNink] (m0) -- (m1);
\draw[->,SNink] (m1) -- (m2);
\draw[->,SNink] (m2) edge[bend left=40] (m0);
\end{tikzpicture}
\par{\fontsize{7.5}{9}\selectfont 循环链接含环，链式存储不能确定树关系}};\pgfextracty{\SNsideheight}{\pgfpointdiff{\pgfpointanchor{SNsidebox}{south}}{\pgfpointanchor{SNsidebox}{north}}}\typeout{SNHEIGHT:6:1:\the\SNsideheight}}
\expandafter\gdef\csname SNside7\endcsname{\node[anchor=north west,inner sep=0pt,text width=48mm,font=\SideFont,align=left] (SNsidebox) at ([xshift=138mm,yshift=-95.00pt]current page.north west) {\textcolor{SNink}{\textbf{最小反例 · 重复申请不等于同时占用}}\par\vspace{1.5mm}循环两次，每次申请一个单元并在下一轮前释放，峰值仍只占一个单元。空间分析要沿时间轴看尚未释放的状态；不能照搬``轮数\ensuremath{\times}每轮空间''的时间计数方法。};\pgfextracty{\SNsideheight}{\pgfpointdiff{\pgfpointanchor{SNsidebox}{south}}{\pgfpointanchor{SNsidebox}{north}}}\typeout{SNHEIGHT:7:9:\the\SNsideheight}}
\expandafter\gdef\csname SNside8\endcsname{\node[anchor=north west,inner sep=0pt,text width=48mm,font=\SideFont,align=left] (SNsidebox) at ([xshift=138mm,yshift=-95.00pt]current page.north west) {\textcolor{SNink}{\textbf{错项排除 · 效率不是唯一质量}}\par\vspace{1.5mm}一个输出错误的程序即使时间和空间都很少，也没有解决目标问题。遇到用效率``全面衡量优劣''的说法，应先检查正确性等前提；正确性、可读性、健壮性不能被遗漏。};\pgfextracty{\SNsideheight}{\pgfpointdiff{\pgfpointanchor{SNsidebox}{south}}{\pgfpointanchor{SNsidebox}{north}}}\typeout{SNHEIGHT:8:11:\the\SNsideheight}}
\expandafter\gdef\csname SNside9\endcsname{\node[anchor=north west,inner sep=0pt,text width=48mm,font=\SideFont,align=left] (SNsidebox) at ([xshift=138mm,yshift=-95.00pt]current page.north west) {\textcolor{SNink}{\textbf{自查提示 · 推导与结果一起核对}}\par\vspace{1.5mm}核对复杂度答案时，同时复查每轮计数、递归终止条件及同时活跃的调用层数。};\pgfextracty{\SNsideheight}{\pgfpointdiff{\pgfpointanchor{SNsidebox}{south}}{\pgfpointanchor{SNsidebox}{north}}}\typeout{SNHEIGHT:9:16:\the\SNsideheight}}
\AddToShipoutPictureFG{\begin{tikzpicture}[remember picture,overlay]\ifcsname SNside\thepage\endcsname\csname SNside\thepage\endcsname\fi\end{tikzpicture}}